\documentclass[11pt]{amsart}
\usepackage{amsmath,amssymb,amsthm}
\usepackage[margin=1.05in]{geometry}

\numberwithin{equation}{section}
\newtheorem{theorem}{Theorem}[section]
\newtheorem{proposition}[theorem]{Proposition}
\newtheorem{lemma}[theorem]{Lemma}
\newtheorem{corollary}[theorem]{Corollary}
\theoremstyle{definition}
\newtheorem{definition}[theorem]{Definition}
\theoremstyle{remark}
\newtheorem{remark}[theorem]{Remark}
\newtheorem{example}[theorem]{Example}

\newcommand{\Z}{\mathbb{Z}}
\newcommand{\Q}{\mathbb{Q}}
\newcommand{\R}{\mathbb{R}}
\newcommand{\N}{\mathbb{N}}
\newcommand{\C}{\mathbb{C}}
\newcommand{\cA}{\mathcal{A}}
\newcommand{\cG}{\mathcal{G}}
\newcommand{\cK}{\mathcal{K}}
\newcommand{\cO}{\mathcal{O}}
\newcommand{\pp}{\mathfrak{p}}
\newcommand{\qq}{\mathfrak{q}}
\newcommand{\aaa}{\mathfrak{a}}
\newcommand{\bb}{\mathfrak{b}}
\newcommand{\Ns}{\mathrm{N}}
\newcommand{\Cl}{\operatorname{Cl}}
\newcommand{\Prob}{\operatorname{Prob}}
\newcommand{\spn}{\operatorname{span}}
\newcommand{\id}{\mathrm{id}}

\PassOptionsToPackage{hyphens}{url}
\usepackage[hidelinks,bookmarks=false]{hyperref}

\title[The fixed-point algebra is the Cartan subalgebra]{Equivariant rigidity of localized Bost--Connes\\ $C^*$-dynamical systems: the fixed-point algebra\\ is the Cartan subalgebra}
\author{Ruqing Chen}
\address{GUT Geoservice Inc., Montreal, Canada}
\email{ruqing@hotmail.com}
\thanks{Sources, code and data for the entire series: \url{https://github.com/Ruqing1963/localized-bost-connes}; this paper is in \texttt{papers/XIII-equivariant/}.}
\date{}

\begin{document}

\begin{abstract}
By Paper IX of the series the von Neumann algebra of a localized Bost--Connes system
remembers nothing: it is injective, hence determined by its flow of weights. There it was
also observed that the obstruction group $\Xi_\beta$ survives at the $C^*$-level, through
the $\mathrm{KMS}_\beta$ simplex. We make that precise and determine how much of the
arithmetic an equivariant isomorphism of $C^*$-dynamical systems recovers.

The technical point on which everything rests is elementary and, as far as we can tell, has
not been isolated. Renault's reconstruction theorem applies to Cartan \emph{pairs}, whereas
an isomorphism of $C^*$-dynamical systems is only required to intertwine the flows; Cartan
subalgebras of a $C^*$-algebra are in general far from unique, so the passage from one to the
other is not automatic. We show it is automatic here, for a reason internal to the system:
\[
\cA_{K,S}^{\sigma}=\overline{\spn}\bigl\{\mu_\aaa f\mu_\bb^*:\Ns\aaa=\Ns\bb\bigr\}
=C(Y_S)\qquad\text{whenever }\Ns:J_S\to\N\text{ is injective,}
\]
since $\Ns\aaa=\Ns\bb$ then forces $\aaa=\bb$ and $\mu_\aaa f\mu_\aaa^*=\alpha_\aaa(f)\in C(Y_S)$.
The Cartan subalgebra is therefore \emph{canonically the fixed-point algebra of the flow},
and every $\sigma$-equivariant isomorphism preserves it. Norm-separation is the same thing
as $\Q$-linear independence of the $\log\Ns\pp$, holds exactly when $S$ contains at most one
prime above each rational prime, and is automatic over $\Q$.

With that, Renault's theorem gives: a $\sigma$-equivariant isomorphism
$\cA_{K_1,S_1}\cong\cA_{K_2,S_2}$ induces an isomorphism of the groupoids of the two systems
intertwining the scaling cocycles $c(\aaa)=\log\Ns\aaa$. The full bisections on which the
cocycle is constant are exactly the semigroup elements, so the isomorphism induces a
norm-preserving bijection $S_1\to S_2$: the multiset $\{\Ns\pp\}_{\pp\in S}$ is recovered. The
symmetry group $G_S$ acts by automorphisms of the groupoid fixing every semigroup
bisection, and over $\Q$, for sets of primes that are Dirichlet-dense at each of their
members, these are the only such automorphisms; whenever the isomorphism intertwines the
symmetry actions --- automatic in that case --- it matches the Frobenius classes and
transports the whole chain $\{\Xi_\beta\}$. In general what is recovered is the compact
space $G_S$, the $\mathrm{KMS}_\beta$ simplices, the transition locus and $\beta_c$; whether the
group structure of $G_S$ is an invariant of the dynamical system for fields with infinite
unit group we do not know.

We are careful about what is \emph{not} recovered. The system sees only the triple
$\bigl(S,\Ns|_S,G_S\bigr)$ with its Frobenius data; the field $K$ enters only through that
triple, and whether the triple determines $K$ is an arithmetic question we do not settle and
which, for $S=P_K$, belongs to the literature on Bost--Connes reconstruction. The narrow class
number is not recovered by any argument here. The KMS bundle is a weaker invariant than the
groupoid: it sees the simplices and their symmetry, but we see no way to extract the
individual norms from it. Finally, the proposed equivariant classification of the Kirchberg
boundary by $KK_\R$ is not available to us: the classification of flows on Kirchberg algebras
up to cocycle conjugacy is not a settled theory, and we do not assert it.
\end{abstract}

\maketitle

\section{The fixed-point algebra}

We use the notation of \cite{Main}: $\cA_{K,S}=C(Y_S)\rtimes J_S$, spanned by $\mu_\aaa f\mu_\bb^*$,
with $\sigma_t(\mu_\aaa)=\Ns\aaa^{it}\mu_\aaa$ and $\alpha_\aaa(f)=\mu_\aaa f\mu_\aaa^*$; it is
the unital full corner of the groupoid algebra of $\cG_S=I_S\ltimes\widetilde Y_S$, and we
write $\cG_S|_{Y_S}$ for the reduction of $\cG_S$ to the compact open subset $Y_S$ of its
unit space, an \'etale groupoid with $C^*(\cG_S|_{Y_S})=\cA_{K,S}$ and $C(Y_S)$ as diagonal.

\emph{The series.} The foundational paper is \cite{Main} and Paper I is \cite{SeriesI};
Papers II, III and IX, cited by numeral, are in preparation, and the facts used from them
are re-proved here.

\begin{definition}\label{def:injective}
Say that $S$ is \emph{norm-separated} if $\Ns:J_S\to\N$, $\aaa\mapsto\Ns\aaa$, is injective.
\end{definition}

\begin{lemma}[When norm-separation holds]\label{lem:sep}
The following are equivalent:
\begin{enumerate}
\item $S$ is norm-separated;
\item the numbers $\log\Ns\pp$, $\pp\in S$, are linearly independent over $\Q$;
\item $S$ contains at most one prime above each rational prime.
\end{enumerate}
In particular every $S$ is norm-separated when $K=\Q$, and $S$ fails to be norm-separated
as soon as it contains two primes above the same rational prime --- two primes of equal
norm, or a prime together with another whose norm is a power of its residue
characteristic.
\end{lemma}

\begin{proof}
(1)$\Leftrightarrow$(2): $J_S$ is free on $S$, so $\Ns\aaa=\Ns\bb$ with $\aaa\ne\bb$ is a
nontrivial relation $\sum_\pp(a_\pp-b_\pp)\log\Ns\pp=0$ with integer coefficients, and every
integer relation arises this way; integer and rational independence agree.
(2)$\Leftrightarrow$(3): $\log\Ns\pp=f_\pp\log p$ for $\pp$ above the rational prime $p$ with
residue degree $f_\pp$, and the $\log p$ are $\Q$-independent by unique factorization; so
a rational relation among the $\log\Ns\pp$ decomposes into relations among the primes above
each fixed $p$, and two primes above the same $p$ satisfy
$f_{\pp'}\log\Ns\pp-f_\pp\log\Ns\pp'=0$. Compare \cite[Lem.~4.1]{SeriesI}.
\end{proof}

\begin{example}\label{ex:fail}
For $K=\Q(i)$: the two primes above $5$ both have norm $5$, so any $S$ containing both fails;
a set containing one prime above each of some rational primes --- one of the two primes
above each $p\equiv1\pmod4$, inert primes $p\equiv3\pmod4$ with norm $p^2$, the prime $1+i$
--- is norm-separated. For $K=\Q(\sqrt2)$ the set of all primes fails at every split prime.
The failure by equal norms is the generic one.
\end{example}

\begin{theorem}[The Cartan subalgebra is the fixed-point algebra]\label{thm:fixed}
If $S$ is norm-separated then
\begin{equation}\label{eq:fixed}
\cA_{K,S}^{\sigma}:=\bigl\{x\in\cA_{K,S}:\sigma_t(x)=x\ \forall t\bigr\}=C(Y_S).
\end{equation}
Consequently the Cartan subalgebra $C(Y_S)\subset\cA_{K,S}$ is canonically determined by
the flow, and every $\sigma$-equivariant isomorphism
$\Phi:\cA_{K_1,S_1}\to\cA_{K_2,S_2}$ satisfies $\Phi\bigl(C(Y_{S_1})\bigr)=C(Y_{S_2})$, hence
is an isomorphism of Cartan pairs.
\end{theorem}

\begin{proof}
A spectral subspace argument: $\sigma_t(\mu_\aaa f\mu_\bb^*)=(\Ns\aaa/\Ns\bb)^{it}\mu_\aaa f\mu_\bb^*$,
so the elements $\mu_\aaa f\mu_\bb^*$ are eigenvectors and $\sigma$ is the restriction to
$\R$ of an action of the compact dual of the countable group $\Gamma_\beta=\langle\Ns\pp\rangle\le\R_{>0}$,
in which $\R$ is dense since $\Gamma_\beta$ is torsion free; averaging over that compact group
is a conditional expectation onto the fixed-point algebra which kills every eigenvector of
nonzero weight, so the fixed-point algebra is the closed span of those $\mu_\aaa f\mu_\bb^*$
with $\Ns\aaa=\Ns\bb$. Norm-separation gives $\aaa=\bb$, and then
$\mu_\aaa f\mu_\aaa^*=\alpha_\aaa(f)$, a continuous function on $Y_S$ supported on the clopen
set $\aaa Y_S$. Hence $\cA^\sigma\subseteq C(Y_S)$; the reverse inclusion holds because
$\sigma_t$ fixes $C(Y_S)$ pointwise. For the consequence, $\Phi\circ\sigma^1_t=\sigma^2_t\circ\Phi$
gives $\Phi(\cA_1^{\sigma^1})=\cA_2^{\sigma^2}$.
\end{proof}

\begin{remark}[Why this needs saying]\label{rem:whyneeded}
Renault's reconstruction theorem is a statement about Cartan \emph{pairs}
$\bigl(B\subset A\bigr)$: it recovers the twisted groupoid from the pair, not from $A$ alone.
A $C^*$-algebra may carry many non-conjugate Cartan subalgebras, so an isomorphism of algebras
--- even an equivariant one --- carries no information about groupoids unless it is known to
respect the pair. Theorem~\ref{thm:fixed} supplies exactly that, by identifying the Cartan
subalgebra intrinsically, in terms of the flow that the isomorphism is assumed to intertwine.
Without it the reconstruction argument of \S\ref{sec:rec} has a gap.
\end{remark}

\begin{remark}[The unseparated case]\label{rem:unsep}
If $S$ is not norm-separated then $\cA^\sigma$ is strictly larger than $C(Y_S)$: it contains
$\mu_\aaa f\mu_{\bb}^*$ for distinct $\aaa,\bb$ of equal norm, which is not in the diagonal.
The fixed-point algebra is then non-commutative and Theorem~\ref{thm:fixed} fails. This does not
show that the reconstruction fails, only that this route to it does; a different intrinsic
characterization of $C(Y_S)$ would be needed, and we do not have one.
\end{remark}

\section{Reconstruction}\label{sec:rec}

Throughout this section $S$ is norm-separated. We first record what in the groupoid
$\cG_S|_{Y_S}$ is intrinsic. Its arrows are the pairs $(\aaa\bb^{-1},y)$ with $y\in Y_S$ and
$\aaa\bb^{-1}y\in Y_S$, the scaling cocycle is $c(\aaa\bb^{-1},y)=\log\Ns\aaa-\log\Ns\bb$,
and for $\aaa\in J_S$ the set $B_\aaa=\{(\aaa,y):y\in Y_S\}$ is a full bisection --- an
open set of arrows on which the source map is a homeomorphism onto the whole unit space
--- with range $\aaa Y_S$ and $c\equiv\log\Ns\aaa$.

\begin{lemma}[The semigroup is intrinsic]\label{lem:bisections}
The full bisections of $\cG_S|_{Y_S}$ on which $c$ is constant are exactly the $B_\aaa$,
$\aaa\in J_S$, and $B_\aaa B_\bb=B_{\aaa\bb}$. Consequently an isomorphism of \'etale
groupoids $\theta:\cG_{S_1}|_{Y_{S_1}}\to\cG_{S_2}|_{Y_{S_2}}$ with $c_2\circ\theta=c_1$ induces a
monoid isomorphism $\pi:J_{S_1}\to J_{S_2}$ with $\Ns\pi(\aaa)=\Ns\aaa$, hence a bijection
$S_1\to S_2$ preserving norms.
\end{lemma}

\begin{proof}
By Lemma~\ref{lem:sep}, norm-separation means that $\Ns:I_S\to\Q_{>0}$ is injective on the
group of fractions $I_S=\Z^{(S)}$, so the value $c(\gamma)$ determines the element
$g\in I_S$ with $\gamma=(g,y)$. A full bisection contains for each $y\in Y_S$ exactly one
arrow with source $y$; if $c$ is constant on it, all these arrows have the same $g=\aaa\bb^{-1}$
with $\aaa,\bb\in J_S$ coprime, and the arrow $(\aaa\bb^{-1},y)$ exists in $\cG_S|_{Y_S}$ only
for $y\in\bb Y_S$. Fullness forces $\bb Y_S=Y_S$, i.e.\ $\bb=1$, so the bisection is $B_\aaa$.
The composition rule is $(\aaa,\bb y)(\bb,y)=(\aaa\bb,y)$. An isomorphism $\theta$ preserving $c$ maps full
bisections with constant $c$ to full bisections with the same constant, and preserves
composition, so $B_\aaa\mapsto B_{\pi(\aaa)}$ defines a monoid isomorphism of the free
abelian monoids $J_{S_i}$ with $\Ns\pi(\aaa)=\Ns\aaa$; a monoid isomorphism of free abelian
monoids maps the free generators to the free generators.
\end{proof}

\begin{definition}\label{def:frobdense}
Let $\operatorname{Aut}_0(\cG_S|_{Y_S})$ be the group of automorphisms of the topological
groupoid $\cG_S|_{Y_S}$ that map each bisection $B_\aaa$ onto itself, with the topology of
uniform convergence on $Y_S$; equivalently, the automorphisms of $\cA_{K,S}$ fixing every
$\mu_\aaa$ and preserving $C(Y_S)$. The translation action of $G_S$ on
$Y_S=\widehat{\cO}_S\times_{U_S}G_S$ commutes with the partial action of $J_S$, so it gives an
injective continuous homomorphism $G_S\to\operatorname{Aut}_0(\cG_S|_{Y_S})$. For $K=\Q$ say
that $S$ is \emph{Dirichlet-dense} if for every $p\in S$ the primes of $S\setminus\{p\}$
generate a dense subgroup of $\Z_p^\times$, i.e.\ meet every residue class modulo every
power of $p$; the set of all primes, and any set of density one in every arithmetic
progression, is Dirichlet-dense.
\end{definition}

\begin{proposition}[The symmetry group is intrinsic over $\Q$]\label{prop:aut}
Let $K=\Q$ and let $S$ be Dirichlet-dense. Then $G_S\to\operatorname{Aut}_0(\cG_S|_{Y_S})$ is an
isomorphism.
\end{proposition}

\begin{proof}
For $K=\Q$, $Y_S=\prod_{p\in S}\Z_p$ with $\aaa\in J_S$ acting by multiplication and
$G_S=\prod_{p\in S}\Z_p^\times$ acting by multiplication as well. Let $\psi\in\operatorname{Aut}_0$.
It commutes with multiplication by every $n\in J_S$ and preserves the ranges $nY_S$, hence
the valuations, hence the set $G_S=Y_S^\times$ of units; write $\psi(u)=u\,c(u)$ for
$u\in G_S$, with $c:G_S\to G_S$ continuous. For $x\in Y_S$ with no zero component, $x=nu$
uniquely with $n\in J_S$, $u\in G_S$, and $\psi(x)=n\psi(u)=x\,c(u)$. Let $p\in S$ and
consider points $x$ with $x_p=0$: they are limits of $n_ku_k$ with $n_k\in J_S$ divisible by
$p^k$, and the $p$-components $u_{k,p}$ of the units are unconstrained, since
$(n_k)_p\to0$; so continuity of $\psi$ at such points forces the components $c(u)_q$, $q\ne p$,
not to depend on $u_p$. Varying $p$ over $S\setminus\{q\}$, the component $c(u)_q$ depends
only on $u_q$: $c(u)=(c_q(u_q))_q$ with $c_q:\Z_q^\times\to\Z_q^\times$ continuous. Now
compare $\psi(nx)$ and $n\psi(x)$ for $x$ a unit and $n\in J_S$: the $q$-component of $nx$ is
$q^{v_q(n)}n'_qx_q$ with $n'_q=n/q^{v_q(n)}$, so $\psi(nx)_q=q^{v_q(n)}n'_qx_q\,c_q(n'_qx_q)$,
while $(n\psi(x))_q=q^{v_q(n)}n'_qx_q\,c_q(x_q)$. Hence $c_q(n'_qx_q)=c_q(x_q)$ for all
$n\in J_S$ and $x_q\in\Z_q^\times$: $c_q$ is invariant under multiplication by the unit parts
of the elements of $J_S$, which generate the subgroup of $\Z_q^\times$ generated by the
primes of $S\setminus\{q\}$, dense by hypothesis. So each $c_q$ is constant, $c$ is a constant
$g\in G_S$, and $\psi$ is multiplication by $g$.
\end{proof}

\begin{remark}\label{rem:autgeneral}
For a general number field the same argument gives less. The constraint from continuity
at the points with infinite valuation at $\pp$ concerns $\psi$ only modulo the image
$r(\cO_\pp^\times)\subseteq G_S$ of the local units, and for fields with infinite unit group
the images of the local units at different primes are not independent in $G_S$, the global
units sitting diagonally in all of them. We do not know whether
$\operatorname{Aut}_0(\cG_S|_{Y_S})=G_S$ in general; the compact model $\overline{\N}^{\,S}\times G_S$
used elsewhere in the series differs from $Y_S$ exactly at these points, and would give the
equality for every Frobenius-dense $S$, but it is not the space on which $\cA_{K,S}$ is built.
\end{remark}

\begin{theorem}[Equivariant reconstruction]\label{thm:rec}
Let $S_1,S_2$ be infinite and norm-separated, and suppose
$\Phi:\bigl(\cA_{K_1,S_1},\sigma^1\bigr)\to\bigl(\cA_{K_2,S_2},\sigma^2\bigr)$ is an
isomorphism of $C^*$-dynamical systems. Then:
\begin{enumerate}
\item $\Phi$ is an isomorphism of Cartan pairs and, by Renault's theorem, is induced by an
isomorphism of \'etale groupoids $\theta:\cG_{S_1}|_{Y_{S_1}}\to\cG_{S_2}|_{Y_{S_2}}$;
\item $\theta$ intertwines the scaling cocycles, $c_2\circ\theta=c_1$;
\item $\theta$ induces a norm-preserving bijection $\pi:S_1\to S_2$, so the multisets
$\{\Ns\pp\}_{\pp\in S_1}$ and $\{\Ns\pp\}_{\pp\in S_2}$ coincide;
\item $\theta$ maps the valuation fibres of $Y_{S_1}$ homeomorphically onto those of
$Y_{S_2}$, compatibly with $\pi$; in particular the fibre over $0$, which is the compact
space $G_{S_1}$, is carried homeomorphically onto $G_{S_2}$;
\item $\Phi$ induces affine isomorphisms $\mathrm{KMS}_\beta(\cA_{K_1,S_1})\cong\mathrm{KMS}_\beta(\cA_{K_2,S_2})$
for every $\beta$, so the compact spaces $G_{S_1}/\Xi_\beta(K_1,S_1)^\perp$ and
$G_{S_2}/\Xi_\beta(K_2,S_2)^\perp$ are homeomorphic, $|\Xi_\beta|$ agree when finite, the
transition loci agree, and $\beta_c(S_1)=\beta_c(S_2)$;
\item suppose $\Phi$ intertwines the symmetry actions, $\Phi\circ\tau_g=\tau_{\iota(g)}\circ\Phi$
for an isomorphism of topological groups $\iota:G_{S_1}\to G_{S_2}$ --- which is automatic,
with $\iota=\theta(\cdot)\theta^{-1}$ on $\operatorname{Aut}_0$, when $K_1=K_2=\Q$ and the
$S_i$ are Dirichlet-dense, by Proposition~\ref{prop:aut}. Then
$\iota(\sigma_\pp)\in\sigma_{\pi(\pp)}\,r(\cO_{\pi(\pp)}^\times)$ for every $\pp\in S_1$: the
Frobenius classes modulo inertia are matched, and $\hat\iota^{-1}$ maps
$\Xi_\beta(K_2,S_2)$ onto $\Xi_\beta(K_1,S_1)$ for every $\beta>0$.
\end{enumerate}
\end{theorem}

\begin{proof}
(1) Theorem~\ref{thm:fixed} makes $\Phi$ an isomorphism of Cartan pairs. The groupoid
$\cG_S|_{Y_S}$ is \'etale, Hausdorff, second countable, amenable, and topologically
principal: the isotropy at a point with all valuations finite is trivial, since
$\aaa\bb^{-1}y=y$ forces $\aaa=\bb$ on such points, and these points are dense. So
$(C(Y_S)\subset\cA_{K,S})$ is a Cartan pair with trivial Weyl twist, and Renault's theorem
\cite{Renault} gives $\theta$ inducing $\Phi$. (2) Under $\Phi=\theta_*$ the flow is
$\sigma_t(f)(\gamma)=e^{itc(\gamma)}f(\gamma)$ on functions supported on bisections, so
$\Phi\sigma^1_t=\sigma^2_t\Phi$ gives $e^{itc_2(\theta\gamma)}=e^{itc_1(\gamma)}$ for all $t$.
(3) is Lemma~\ref{lem:bisections}. (4) $\theta$ maps $B_\aaa$ to $B_{\pi(\aaa)}$, hence the
ranges $\aaa Y_{S_1}$ to $\pi(\aaa)Y_{S_2}$, hence the valuation fibres to the valuation
fibres; the fibre over $0$ is $U_S\times_{U_S}G_S=G_S$. (5) An isomorphism of
$C^*$-dynamical systems carries $\mathrm{KMS}_\beta$ states to $\mathrm{KMS}_\beta$ states; the
simplices are $\Prob(G_S/\Xi_\beta^\perp)$ by \cite[Thm.~3.6]{Main}, and an affine isomorphism
of Bauer simplices is a homeomorphism of their extreme boundaries; the transition locus is
the set of $\beta$ at which the simplex changes. (6) Since $\Phi=\theta_*$ preserves the
Cartan pair, intertwining the symmetry actions means $\theta\circ\tau_g=\tau_{\iota(g)}\circ\theta$
on $Y_{S_1}$; when $\operatorname{Aut}_0=G_S$ on both sides this holds with
$\iota=\theta(\cdot)\theta^{-1}$, an isomorphism of topological groups since $\theta$ matches
the bisections $B_\aaa$. On the valuation fibres, which are $G_S$-torsors, $\theta$ is
$\iota$-equivariant, so in the coordinates $y=(v,h)$ of \cite[Lem.~2.10]{Main} on the points
with finite valuations, $\theta(v,h)=(\pi(v),\iota(h)c_v)$ with constants $c_v\in G_{S_2}$, and
$\theta(\pp\cdot(v,h))=\pi(\pp)\cdot\theta(v,h)$ reads
$\iota(\sigma_\pp)\iota(h)c_{v+e_\pp}=\sigma_{\pi(\pp)}\iota(h)c_v$, i.e.\
$c_{v+e_\pp}=u_\pp c_v$ with $u_\pp=\iota(\sigma_\pp)^{-1}\sigma_{\pi(\pp)}$. The points
$(ke_\pp,h)$ converge in $Y_{S_1}$ as $k\to\infty$ modulo the inertia image $r(\cO_\pp^\times)$,
and their images $(ke_{\pi(\pp)},\iota(h)u_\pp^kc_0)$ must converge modulo
$r(\cO_{\pi(\pp)}^\times)$; a sequence $u^k$ converging in the compact group
$G_{S_2}/r(\cO_{\pi(\pp)}^\times)$ has $u$ trivial there. So $u_\pp\in r(\cO_{\pi(\pp)}^\times)$.
Finally $\Xi_\beta$ is characterized by the series $\sum_\pp\Ns\pp^{-\beta}|1-\chi(\sigma_\pp)|^2$
\cite[(1.1)]{Main}, whose convergence class is unchanged when $\sigma_\pp$ is altered by
inertia elements at $\pp$ for all but finitely many $\pp$, the $\mathrm{KMS}$ simplex being
an invariant of the system; so it is transported by $\pi$ and $\hat\iota$.
\end{proof}

\begin{remark}[The hypotheses]\label{rem:indep}
Norm-separation enters twice, in Theorem~\ref{thm:fixed} and in
Lemma~\ref{lem:bisections}; by Lemma~\ref{lem:sep} it is the same as the $\Q$-independence
of the $\log\Ns\pp$, so there is only one arithmetic condition on $S$ in parts (1)--(5), and
it is automatic over $\Q$. Part (6) needs the symmetry actions to be intertwined: without
that, $\theta$ still matches the fibres and the simplices, but a homeomorphism of compact
groups carries no group structure, and we do not know whether $G_S$ with its Frobenius
classes is an invariant of the dynamical system in general (Remark~\ref{rem:autgeneral}).
\end{remark}

\section{What is not recovered}

\begin{remark}[The field]\label{rem:field}
Theorem~\ref{thm:rec} recovers the pair $\bigl(S,\Ns|_S\bigr)$, and, when the symmetry
actions are intertwined, the triple $\bigl(S,\Ns|_S,G_S\bigr)$ with its Frobenius data. It does
\emph{not} recover $K$. The system is built from that triple and nothing else, so two
number fields yielding the same triple give isomorphic dynamical systems, whatever their
other invariants. Whether the triple determines $K$ is an arithmetic question, not an
operator-algebraic one; for $S=P_K$ it is the subject of the literature on reconstruction from
Bost--Connes systems, which we have not been able to consult here and against which any claim
of field reconstruction should be checked. We make no such claim.
\end{remark}

\begin{remark}[Arithmetic equivalence is not obviously the right equivalence]\label{rem:arith}
Arithmetically equivalent fields share their Dedekind zeta function, hence the multiset of
norms $\{\Ns\pp\}$ for $S=P_K$. They need not share $G_K$, which involves the class group and
the units, and $G_S$ is part of the recovered data by
Theorem~\ref{thm:rec}(6). So the invariant recovered here is finer than the zeta function and
coarser than $K$; locating it exactly is open. We note that the narrow class number
$h^+_S=|\Cl^+_S|$ is not recovered by any argument of this paper: $\Cl^+_S$ is the quotient of
$G_S$ by the image of the local units, a finite-index open subgroup which nothing in the
dynamical system singles out.
\end{remark}

\section{The KMS bundle is a weaker invariant}

\begin{definition}\label{def:bundle}
The \emph{KMS bundle} of $(\cA_{K,S},\sigma)$ is
$\bigsqcup_{\beta>0}\mathrm{KMS}_\beta(\cA_{K,S})$ over $(0,\infty)$, with the fibrewise
affine structure and the $G_S$-action. An isomorphism of KMS bundles is a fibrewise affine
bijection over the identity of $(0,\infty)$ together with an isomorphism $\iota$ of the
symmetry groups, the two being compatible.
\end{definition}

\begin{proposition}[What the bundle sees]\label{prop:bundle}
A fibrewise affine bijection of KMS bundles over the identity of $(0,\infty)$ determines the
compact spaces $G_S/\Xi_\beta^\perp$ for every $\beta$, hence the transition locus, $\beta_c$,
and $|\Xi_\beta|$ where finite. An isomorphism of KMS bundles in the sense of
Definition~\ref{def:bundle} determines in addition the chain $\{\Xi_\beta\}_{\beta>0}$ as
subgroups of $\widehat{G_S}$, since $\Xi_\beta^\perp$ is the stabilizer in $G_S$ of any extreme
point of the fibre at $\beta$.
\end{proposition}

\begin{proof}
Each of these is read off \cite[Thm.~3.6]{Main}: the fibre at $\beta$ is
$\Prob(G_S/\Xi_\beta^\perp)$, on whose extreme points $G_S$ acts transitively with
stabilizer $\Xi_\beta^\perp$.
\end{proof}

\begin{remark}[What the bundle does not see]\label{rem:bundleweak}
We see no way to extract the individual norms $\Ns\pp$ from the bundle. The fibres depend on
$S$ only through the sets $\{\pp:\chi(\sigma_\pp)\ne1\}$ and their convergence exponents; two
systems whose primes have different norms but the same detecting sets and the same
convergence behaviour would have isomorphic bundles. The bundle is therefore weaker than the
groupoid of Theorem~\ref{thm:rec}, which recovers the norms outright, and the chain of
implications runs
\[
\text{equivariant }C^*\text{-isomorphism}\ \Longrightarrow\ \text{groupoid isomorphism}\ \Longrightarrow\ \text{KMS bundle isomorphism},
\]
with the converses unproved. We do not claim the bundle is a complete invariant.
\end{remark}

\section{The Kirchberg boundary: a caution}

\begin{remark}[Flows on Kirchberg algebras]\label{rem:kirchberg}
By Paper II of the series the boundary $\partial\cA^{\mathrm{aff}}_{K,S}$ of the affine
system is a unital Kirchberg algebra satisfying the UCT, and it carries the flow
$\sigma_t(\mu_{(a,\pp)})=\Ns\pp^{it}\mu_{(a,\pp)}$. It is tempting to classify the pair by
equivariant $KK$. We do not, for two reasons.

First, the classification of \emph{flows} --- actions of $\R$ --- on Kirchberg algebras up to
cocycle conjugacy is not a settled theory. The Kirchberg--Phillips theorem
\cite{KirchbergPhillips} classifies the algebras; the equivariant theory is well developed
for finite and, more recently, for many amenable \emph{discrete} groups \cite{Szabo}, but
$\R$-actions on purely infinite algebras remain substantially open, and we are not able to
invoke a classification theorem here. (For the diagonal-preserving isomorphisms of
\S\ref{sec:rec}, by contrast, the relevant rigidity is that of continuous orbit
equivalence \cite{Li,Matsumoto}, whose groupoid form is Renault's theorem.)

Second, and more concretely, by Paper II this system has
$\mathrm{KMS}_\beta$ states only at $\beta=1$. Whatever the classification of the flow, the
equilibrium invariant that carries the arithmetic in the Toeplitz system has collapsed to a
single point on the boundary, so the boundary is not where the reconstruction of
\S\ref{sec:rec} can be carried out.
\end{remark}

\section{Assessment}

\[
\begin{array}{@{}l@{\ \ }l@{\ \ }l@{}}
\text{Statement} & \text{Status} & \text{Reference}\\\hline
\cA^\sigma=C(Y_S)\ \text{for norm-separated }S & \textbf{true} & \text{Thm.~\ref{thm:fixed}}\\
\text{norm-separation automatic over }\Q & \text{true} & \text{Lem.~\ref{lem:sep}}\\
\text{equivariant iso preserves the Cartan pair} & \textbf{true} & \text{Thm.~\ref{thm:fixed}}\\
\text{Renault: pair}\Rightarrow\text{groupoid} & \text{true} & \text{Thm.~\ref{thm:rec}(1)}\\
\text{norm-separation}=\Q\text{-indep.\ of }\log\Ns\pp & \text{true} & \text{Lem.~\ref{lem:sep}}\\
\text{norms }\{\Ns\pp\}\ \text{recovered} & \text{true} & \text{Thm.~\ref{thm:rec}(3)}\\
\text{simplices, transition locus, }\beta_c\ \text{recovered} & \text{true} & \text{Thm.~\ref{thm:rec}(5)}\\
G_S,\ \sigma_\pp,\ \{\Xi_\beta\}\ \text{recovered} & \text{if symmetries intertwined} & \text{Thm.~\ref{thm:rec}(6)}\\
\operatorname{Aut}_0=G_S & \text{true over }\Q\text{, Dirichlet-dense }S & \text{Prop.~\ref{prop:aut}}\\
h^+_S\ \text{recovered} & \textbf{not claimed} & \text{Rem.~\ref{rem:arith}}\\
\text{the field }K\ \text{recovered} & \textbf{not claimed} & \text{Rem.~\ref{rem:field}}\\
\text{KMS bundle a complete invariant} & \textbf{not claimed; weaker} & \text{Rem.~\ref{rem:bundleweak}}\\
\text{Kirchberg boundary classified by }KK_\R & \textbf{not available} & \text{Rem.~\ref{rem:kirchberg}}
\end{array}
\]

The contrast with Paper IX is the point. There the von Neumann algebra was shown to
retain nothing, the groupoid being amenable and Connes' classification leaving only the flow
of weights. Here the $C^*$-algebra \emph{with its flow} retains the norms, the
$\mathrm{KMS}$ simplices, and --- once the symmetries are intertwined, which over $\Q$ is
automatic --- the symmetry group with its Frobenius classes and the whole chain
$\{\Xi_\beta\}$. The mechanism is
Theorem~\ref{thm:fixed}: the flow itself pins down the Cartan subalgebra, so the
reconstruction theory of Renault becomes available to a merely equivariant isomorphism. Weak
closure destroys the flow --- the modular group of a single extremal state is an inner
datum --- and with it the pin.

Four questions. Is $C(Y_S)$ intrinsically characterized when $S$ is not norm-separated, so
that Theorem~\ref{thm:rec} extends beyond Lemma~\ref{lem:sep}? Is
$\operatorname{Aut}_0(\cG_S|_{Y_S})=G_S$ for fields with infinite unit group, so that the group
structure of $G_S$ is an invariant of the dynamical system there too? Is the KMS bundle a
complete invariant after
all, or is there a pair of systems with isomorphic bundles and non-isomorphic groupoids ---
we would expect the latter and suggest looking at sets with the same detecting data but
different norms. And what exactly is the arithmetic content of the recovered triple
$(S,\Ns|_S,G_S)$: it is finer than the Dedekind zeta function, since $G_S$ carries class group
and unit data, and coarser than $K$; where it sits between them we do not know.

\begin{thebibliography}{9}
\bibitem{KirchbergPhillips} E. Kirchberg, N. C. Phillips, \emph{Embedding of exact $C^*$-algebras in the Cuntz algebra $\mathcal O_2$}, J. Reine Angew. Math. 525 (2000), 17--53.
\bibitem{Li} X. Li, \emph{Continuous orbit equivalence rigidity}, Ergodic Theory Dynam. Systems 38 (2018), 1543--1563.
\bibitem{Matsumoto} K. Matsumoto, H. Matui, \emph{Continuous orbit equivalence of topological Markov shifts and Cuntz--Krieger algebras}, Kyoto J. Math. 54 (2014), 863--877.
\bibitem{Renault} J. Renault, \emph{Cartan subalgebras in $C^*$-algebras}, Irish Math. Soc. Bull. 61 (2008), 29--63.
\bibitem{Szabo} G. Szab\'o, \emph{On the classification of $C^*$-dynamical systems}, survey, 2021.
\bibitem{Main} R. Chen, \emph{KMS state uniqueness and phase transitions in localized Bost--Connes systems: from Chebotarev density to the Hardy--Littlewood conjecture}, preprint, 2026, DOI \href{https://doi.org/10.5281/zenodo.22152101}{10.5281/zenodo.22152101}.
\bibitem{SeriesI} R. Chen, \emph{Localized Bost--Connes systems I: factor type and the spectrum of transitions}, preprint, 2026, DOI \href{https://doi.org/10.5281/zenodo.22160827}{10.5281/zenodo.22160827}.
\end{thebibliography}

\end{document}
